%%%%%%%%%%%%%%%%%%%%%%%%%%%%%%%%%%%%%%%%%%%%%%%%%%%%%%%
% SECTION 7: RESULTS 2.4 - NATURAL LANGUAGE QUERY-TO-SQL INTERFACE
%%%%%%%%%%%%%%%%%%%%%%%%%%%%%%%%%%%%%%%%%%%%%%%%%%%%%%%
% 
% INPUTS NEEDED: 
%   - Training data generation details
%   - Model accuracy metrics
%   - Baseline comparisons
%   - Latency measurements
%   - Error analysis
%
% FIGURES: Figure 2 (figure2.png) - NL-to-SQL system architecture
% TABLES: None
% DEPENDENCIES: Methods section, Supplementary Methods
%
%%%%%%%%%%%%%%%%%%%%%%%%%%%%%%%%%%%%%%%%%%%%%%%%%%%%%%%

\subsection{Natural language query-to-SQL interface performance}\label{subsec4}

%%% PARAGRAPH 1: Problem statement - SQL barrier %%%
Despite successful automated evidence extraction, clinical adoption of genomic databases remains limited by the requirement for structured query languages that are inaccessible to most clinicians. Current database interfaces require users to understand complex schema relationships and SQL syntax, creating a significant barrier between curated genomic evidence and clinical decision-making. To address this accessibility challenge, we developed a lightweight, domain-specific access layer that enables natural language querying of the curated OncoCITE-KB through automated SQL translation.

%%% PARAGRAPH 2: Training data and model training → REFERENCES FIGURE 2 %%%
Automated template expansion produced 21,000 question--SQL pairs distributed across four complexity tiers, enabling comprehensive evaluation of system performance across varying query sophistication levels. The LoRA-augmented Llama-3 3B model converged after three epochs, with training loss stabilization indicating adequate pattern capture without overfitting as shown in Fig.~\ref{fig2}.

%%% PARAGRAPH 3: Performance results %%%
When evaluated on the complete 21,000 synthetic question-query test set, the domain-tuned Llama-3 3B model achieved 83.24\% accuracy in generating execution-correct SQL queries. Generation latency remained between 6.0 and 7.4 seconds across all complexity tiers, indicating effective internalization of the CIViC database schema for handling intricate joins and aggregations with consistent performance regardless of query complexity.

%%% PARAGRAPH 4: Baseline comparison and error analysis %%%
Comparative evaluation against baseline systems demonstrated the advantage of domain-specific fine-tuning. A general Civic-SQL-Large model achieved only 55\% overall accuracy with sharp performance degradation as complexity increased, reaching only 40\% for complex queries and 20\% for super-complex queries. A separately fine-tuned Llama-SQL variant performed comparably on simple requests with approximately 50\% overall accuracy yet returned no correct output for super-complex queries. In both baseline systems, 84\% of errors stemmed from schema misalignment, producing syntactically valid SQL that referenced non-existent tables or columns. This module enables access to OncoCITE-KB. The knowledge-management component (Phase~3) described in the next section provides clinical interpretation.

%%%%%%%%%%%%%%%%%%%%%%%%%%%%%%%%%%%%%%%%%%%%%%%%%%%%%%%
% FIGURE 2 REQUIREMENTS (figure2.png):
%
% TITLE: "Natural language to SQL query system architecture"
%
% MUST SHOW:
%   - Template expansion process → 21K pairs
%   - Four complexity tiers with distribution
%   - LoRA fine-tuning on Llama-3 3B
%   - Training convergence (3 epochs)
%   - Query processing flow
%
% PERFORMANCE CALLOUTS TO INCLUDE:
%   - 83.24% execution accuracy
%   - 6-7 second latency
%   - 28-point improvement over baseline (55%)
%   - Four tiers: Simple 5K, Medium 5K, Complex 5K, Super-Complex 6K
%
% CAPTION TEXT:
% "Natural language to SQL query system architecture. Domain-specific 
% text-to-SQL translation using LoRA fine-tuning on synthetic training 
% data. The system achieves 83.24% execution accuracy across four 
% complexity tiers (Simple: 5K pairs single-table; Medium: 5K pairs 
% multi-join; Complex: 5K pairs aggregations; Super-Complex: 6K pairs 
% nested subqueries) with consistent 6--7 second latency, representing 
% a 28-point improvement over baseline (55%)."
%
%%%%%%%%%%%%%%%%%%%%%%%%%%%%%%%%%%%%%%%%%%%%%%%%%%%%%%%

%%%%%%%%%%%%%%%%%%%%%%%%%%%%%%%%%%%%%%%%%%%%%%%%%%%%%%%
% KEY METRICS IN THIS SECTION:
%
% TRAINING DATA:
%   □ 21,000 question-SQL pairs
%   □ Four complexity tiers
%   □ 3 epochs to convergence
%
% MODEL:
%   □ Llama-3 3B with LoRA
%   □ Domain-specific fine-tuning
%
% PERFORMANCE:
%   □ 83.24% execution-correct accuracy
%   □ 6.0-7.4 second latency (consistent across tiers)
%
% BASELINE COMPARISONS:
%   ┌──────────────────────┬─────────┬─────────┬──────────────┐
%   │ Model                │ Overall │ Complex │ Super-Complex│
%   ├──────────────────────┼─────────┼─────────┼──────────────┤
%   │ OncoCITE (Llama-3)   │ 83.24%  │   --    │      --      │
%   │ Civic-SQL-Large      │ 55%     │  40%    │     20%      │
%   │ Llama-SQL variant    │ ~50%    │   --    │      0%      │
%   └──────────────────────┴─────────┴─────────┴──────────────┘
%
% ERROR ANALYSIS:
%   □ 84% of baseline errors = schema misalignment
%   □ Syntactically valid SQL referencing non-existent tables/columns
%
% COMPLEXITY TIER DISTRIBUTION:
%   □ Simple: 5K pairs (single-table queries)
%   □ Medium: 5K pairs (2-3 table joins)
%   □ Complex: 5K pairs (aggregations)
%   □ Super-Complex: 6K pairs (nested subqueries)
%
%%%%%%%%%%%%%%%%%%%%%%%%%%%%%%%%%%%%%%%%%%%%%%%%%%%%%%%

%%%%%%%%%%%%%%%%%%%%%%%%%%%%%%%%%%%%%%%%%%%%%%%%%%%%%%%
% TRANSITION NOTE:
% This section ends with: "This module enables access to OncoCITE-KB. 
% The knowledge-management component (Phase 3) described in the next 
% section provides clinical interpretation."
%
% This bridges to Results 2.5 (Knowledge-Enhanced Interface)
%%%%%%%%%%%%%%%%%%%%%%%%%%%%%%%%%%%%%%%%%%%%%%%%%%%%%%%
